\section{Connecting sequence processing to the database}
\label{sec:arch-core-boundaries}

Consider a search between stored reference sequences and query sequences.
Both tables below have an identifier \texttt{id} and a sequence column
\texttt{seq}. One statement supplies comparisons to the alignment system:

\noindent\begin{minipage}{\linewidth}
\begin{lstlisting}[style=bioseq-arch-sql,
  caption={A sequence join connects SQL input to the processing core.},
  label={lst:arch-opening-query}]
SELECT r.id AS target_id, q.id AS query_id
FROM refs AS r JOIN queries AS q
  ON mmsearch(r.seq, q.seq);
\end{lstlisting}
\end{minipage}

PostgreSQL supplies the rows and owns their visibility. BioSeqDB selects
compatible implementations, obtains sequence and index views, schedules the
comparisons, and returns matches for SQL projection. This chapter explains
how those responsibilities fit together.

A \emph{batch} groups input comparisons for processing. A \emph{wave} is a
subset admitted together under the memory policy. A \emph{task} is a unit of
work submitted to the task scheduler and can process one or more comparisons.
These terms describe how work is grouped, not different alignment algorithms.

The preceding chapters describe the work performed by a sequence search.
Minimizers produce sparse seed keys
(Section~\ref{sec:minimizers-role}), chaining groups matching positions
(Section~\ref{sec:chaining-bioseqdb}), and final alignment compares the selected
intervals at base level (Section~\ref{sec:final-alignment-role}). These stages
must also fit into a database execution. The sequence store determines how
symbols can be read, persistent indexes retain lookup work, and the custom
scan turns matches into SQL rows.

BioSeqDB connects these components through a reusable C++ core and
PostgreSQL-specific adapters. The core expresses biological operations over
logical sequences and index views. The adapters connect those operations to
database resources. This division gives three kinds of code different duties:
\begin{itemize}
  \item \textbf{Algorithms} consume symbols, positions and scores. Their
    interfaces do not require a PostgreSQL tuple or relation.
  \item \textbf{Adapters} construct suitable views and keep their backing
    resources valid. They also convert returned results into database values.
  \item \textbf{PostgreSQL execution} controls planning, snapshots, locks and
    the lifetime of a statement. These mechanisms remain native to PostgreSQL.
\end{itemize}

The build reflects this division. A shared object-library target owns the
PostgreSQL-free modules, and the extension consumes their compiled objects
and module interfaces. PostgreSQL-dependent modules are compiled separately.
The common code is therefore shared with non-server consumers without
recompiling it under PostgreSQL-specific definitions. This is a source and
compilation boundary. It does not by itself define a separately distributed
library or remove PostgreSQL discovery from the project's build configuration.

\subsection{Reuse at the alphabet boundary}
\label{sec:arch-seqan-reuse}

SeqAn provides an established example of composing efficient sequence-analysis
software from generic C++ components. Its library includes sequence
representations, indexes and configurable alignment algorithms, with SIMD
alignment already described in the 2017 overview.\autocite{REINERT2017157,seqanManual}
BioSeqDB uses SeqAn3's alphabet layer within its own representation. A rank
is a small integer code for a symbol, and the alphabet's cardinality gives
the number of possible ranks. Rank and character conversions
give those ranks their biological meaning, and complement operations apply
where the alphabet supports them.\autocite{seqan3AlphabetsDocs}
The implementation also uses related type traits and a few mathematical
utilities.

BioSeqDB defines the packed storage, lazy views and search composition around
these alphabet semantics. This lets the same logical symbol model participate
in storage reads and in the local alignment kernels. It also lets index
preparation produce the lookup form needed by a particular execution.
Keeping this composition local gives BioSeqDB control over when storage is
read, which stages run, and how long their data survives. SeqAn's own generic
interfaces and optimized algorithms remain a precedent for that design.
The choice of a local pipeline does not establish a performance advantage
over a pipeline assembled from SeqAn components.

Two applications illustrate why composition matters. Masai describes seed
search over conceptual trees but chooses concrete index structures for each
role and combines the search with banded bit-vector
verification.\autocite{siragusa2012fastsensitivereadmapping}
DREAM-Stellar groups query segments by reference region and uses bounded
batches with shared sequence ownership before local
verification.\autocite{Aasna2026}
These are examples of separating algorithms, representations and scheduling.
BioSeqDB makes those choices around its database inputs and result contracts.

\section{Contracts describe the operations a stage needs}
\label{sec:arch-operation-contracts}

The implementation is built in C++26 mode. It uses templates, concepts and
named modules to express these boundaries, although concepts and modules
entered the language in C++20. A concept states which operations a type must
provide. A template constrained by that concept can accept different
representations without requiring inheritance from one common base class.

The following reduced API excerpt shows the contiguous minimizer-lookup
contract. It keeps the current operation and return type, while shortening
the template parameter names. A lookup receives a hash and returns a borrowed
span of encoded positions.

\noindent\begin{minipage}{\linewidth}
\begin{lstlisting}[style=bioseq-cpp26,numbers=none,
  basicstyle=\footnotesize\ttfamily,
  caption={The contiguous lookup API constrains the operation visible to a consumer.},
  label={lst:arch-lookup-contract}]
template<class I>
concept contiguous_minimizer_lookup_index =
    requires(const I& index, std::uint64_t hash) {
        { index.lookup(hash) }
            -> std::same_as<std::span<const std::uint64_t>>;
    };
\end{lstlisting}
\end{minipage}

This is one supported lookup shape. Other consumers can receive segmented
position spans or canonical anchors prepared for a query. The consumer needs
the meaning and lifetime of those positions, rather than the PostgreSQL
relation from which they came. In particular, the total occurrence count
remains authoritative when a selective reader omits the postings of a
high-frequency seed. The storage choice must preserve the pruning decision.

Index construction uses a separate contract. It describes typed chunk reads
and writes, enumeration and temporary sorted-run storage. Concrete
PostgreSQL types satisfy the contracts for per-fragment chunks, a single
per-sequence value, and table-wide chunks. Compile-time assertions check the
required interfaces. Separating these storage operations from the query's
lookup operation lets each side use the representation suited to its work.
The persistence mechanisms are developed in
Section~\ref{sec:index-storage-role}.

Sequence access follows the same principle. A scanner supplies size and
fragment-access operations. Its PostgreSQL implementation reads the fragments
described in Section~\ref{sec:pg-storage-role}, while an in-memory
implementation exposes caller-owned packed
storage. Algorithms normally use the logical views introduced in
Section~\ref{sec:sequence-abstractions}, so they do not need to repeat either
scanner's storage logic. A smaller example appears in column comparison:
a result-sink concept asks for result allocation and writable value and null
masks. The PostgreSQL sink supplies those buffers using database allocation
rules. Character comparison that requires PostgreSQL collation remains in
the PostgreSQL adapter.

These are selective abstractions. A minimizer lookup needs occurrence
positions, but it does not need an abstract join operator. A sequence reader
needs valid source access, but it does not implement an alternative MVCC
system. Keeping the database mechanisms in PostgreSQL preserves their
existing ownership and visibility rules. The common contracts instead expose
the operations needed by the biological stages and their test callers.

\section{Ownership completes the interface}
\label{sec:arch-ownership-boundary}

A successful concept check proves that an expression has a usable type.
It does not prove that a returned span remains valid or that its source may
be read from another thread. Figure~\ref{fig:arch-core-boundary} separates the
operation interface from the resources that make an execution valid.

\begin{figure}[htbp]
\centering
\begin{tikzpicture}[x=1cm,y=1cm,font=\footnotesize,
  arch owner/.style={draw=blue!55!black,fill=blue!5,align=center,
    inner sep=5pt,minimum height=.85cm},
  arch adapter/.style={draw=black!60,fill=black!3,align=center,
    inner sep=5pt,minimum height=.85cm},
  arch flow/.style={-{Latex[length=1.8mm]},thick}]
  \node[arch owner,text width=12.25cm] (owner) at (6.5,5.1)
    {\textbf{PostgreSQL executor owns the resources}\\
     locked relation, registered snapshot, memory contexts};
  \node[arch adapter,text width=5.15cm] (source) at (3.2,3.65)
    {\textbf{Source and index adapters}\\
     keep source data available};
  \node[arch adapter,text width=5.15cm] (projection) at (9.8,3.65)
    {\textbf{Result projection}\\
     construct database values and SQL rows};
  \draw[arch flow] (3.2,4.64)--(source.north);
  \draw[arch flow] (projection.north)--(9.8,4.64);

  \draw[dashed,black!55] (.15,2.65)--(12.85,2.65);
  \node[fill=white,inner sep=4pt] at (6.5,2.65)
    {module interface boundary};
  \draw[draw=green!45!black,fill=green!4]
    (.2,-.18) rectangle (12.8,1.65);
  \node[font=\footnotesize\bfseries] at (6.5,1.28)
    {Sequence-processing core};
  \node at (6.5,.84)
    {minimizers $\rightarrow$ lookup $\rightarrow$ chaining $\rightarrow$ alignment};
  \node[align=center] at (6.5,.32)
    {Worker execution additionally requires shareable input\\
     retained through completion, and its own temporary storage.};
  \draw[arch flow] (source.south)--(3.2,1.65)
    node[pos=.75,left,align=right] {logical views\\and index views};
  \draw[arch flow] (9.8,1.65)--(projection.south)
    node[pos=.25,right,align=left] {owning matches\\and CIGAR};
\end{tikzpicture}
\caption{Data crosses a narrow interface while the PostgreSQL executor
retains database ownership. Arrows show runtime data flow. Static module
dependencies point from the adapters to the core contracts. A logical view
may still request a PostgreSQL read on the owning backend thread. Worker
execution requires an additional check of input sharing and lifetime.}
\label{fig:arch-core-boundary}
\end{figure}

The custom scan from Section~\ref{sec:postgres-custom-scan-scope} owns its
locked relation and registered snapshot. Read
adapters borrow them within the statement's lifetime. A lazy sequence view
can therefore be PostgreSQL-free as a C++ interface while still requesting
data through a PostgreSQL-backed source. The source rejects calls from the
wrong thread before entering PostgreSQL error-handling machinery. The
sharing rules in Section~\ref{sec:sequence-author-contract} apply here:
worker eligibility and source lifetime are separate requirements. Depending
on the route, preparation may copy data into admitted resident storage or
retain an immutable mapped index until its readers finish.

Memory also has several owners. The core context separates temporary
alignment storage from longer-lived pipeline storage. Numerical worker tasks
install a context for their executing thread, while database buffers remain under PostgreSQL memory
contexts. A memory permit accounts for a reservation but does not own the
allocation itself. The bytes must be released before that reservation becomes
available to another task. Returned matches and their CIGAR operations use
owning C++ vectors, which the PostgreSQL adapter can consume without retaining
a temporary alignment arena.

Error handling follows the same division. PostgreSQL failures are translated
at checked calls so that surrounding C++ owners can unwind normally.
Database resource cleanup also has to remain valid on PostgreSQL's own error
path. PostgreSQL interrupt checks occur at supported adapter boundaries,
including between rows, rather than inside every native numerical loop.
Thus a generic function signature alone makes no promise about cancellation
latency.

The remaining shared assumptions are concrete. Both sides must agree on
alphabet ranks, coordinate conventions and index encodings. A prepared index
must describe the source version and minimizer rules expected by the query.
The build must contain the selected implementation. These agreements make
the stages composable, while the adapters keep each execution within the
database's resource and visibility rules.
